\documentclass[12pt,a4paper]{article}
\usepackage[a4paper,left=22mm,right=22mm,top=23mm,bottom=24mm]{geometry}
\usepackage{iftex}
\ifPDFTeX
  \usepackage[T2A]{fontenc}
  \usepackage[utf8]{inputenc}
\else
  \usepackage{fontspec}
  \setmainfont{DejaVu Serif}
\fi
\usepackage[russian]{babel}
\usepackage{amsmath,amssymb,mathtools}
\usepackage{array,booktabs,longtable}
\usepackage{xcolor}
\usepackage{tikz}
\usepackage[colorlinks=true,linkcolor=blue!55!black]{hyperref}
\usepackage{microtype}
\setlength{\parindent}{0pt}
\setlength{\parskip}{6pt}
\setlength{\emergencystretch}{2em}
\renewcommand{\arraystretch}{1.23}
\hypersetup{pdftitle={Ревью домашней работы},pdfauthor={Лёвин Артём Александрович}}
\begin{document}
\thispagestyle{empty}
\vspace*{0.31\textheight}
\begin{center}
{\Huge\bfseries Ревью домашней работы\par}
\vspace{22mm}
{\large Автор: Лёвин Артём Александрович\par}
\vspace{6mm}
{\large 20 октября 2026 года\par}
\end{center}
\vfill
\begin{center}
\begin{tikzpicture}
\draw[blue!30!black, line width=0.65pt] (0,0) .. controls (2.2,0.60) and (4.6,-0.60) .. (6.8,0);
\end{tikzpicture}
\end{center}
\newpage
\section*{Сводная таблица заданий, требующих разбора}
\small
\begin{longtable}{@{}p{0.070\linewidth}p{0.193\linewidth}p{0.271\linewidth}p{0.155\linewidth}p{0.175\linewidth}@{}}
\toprule
\textbf{№} & \textbf{Тема} & \textbf{Суть ошибки} & \textbf{Ваш ответ} & \textbf{Правильный ответ}\\
\midrule
\endfirsthead
\toprule
\textbf{№} & \textbf{Тема} & \textbf{Суть ошибки} & \textbf{Ваш ответ} & \textbf{Правильный ответ}\\
\midrule
\endhead
\hyperref[task:1b]{1б} & Корень третьей степени & Пункт не выполнен: значение корня не найдено. & Не указан & $3{,}5$\\
\addlinespace[4pt]
\hyperref[task:10]{10} & Квадратичная зависимость пути от времени & Разность времён ошибочно возведена в квадрат вместо вычитания двух пройденных расстояний. & $0{,}2$ м & $1$ м\\
\bottomrule
\end{longtable}
\normalsize
\newpage
\thispagestyle{empty}
\vspace*{\fill}
\begin{center}
{\LARGE\bfseries Разбор ошибок}
\end{center}
\vspace*{\fill}
\newpage
\phantomsection\label{task:1b}
\section*{Задание 1б. Корень третьей степени}
\textbf{Текст задания.} Вычислите:
\[
\sqrt[3]{42{,}875}.
\]

\textbf{Ваше решение.} На предоставленных фотографиях есть вычисление пункта 1а, но решение пункта 1б не представлено.

\textbf{Ваш ответ.} Не указан.

\textcolor{red}{\textbf{Ошибка:}} пункт задания остался без решения и ответа. Здесь обнаружен не неверный вычислительный переход, а отсутствие необходимого вычисления.

\textbf{Где именно возникает неполнота.} В задании 1 требуется выполнить два действия: пункт «а» и пункт «б». После вычисления пункта «а» должен быть отдельно рассмотрен корень третьей степени из числа $42{,}875$. На фотографиях этого шага нет; поэтому последнего правильного и первого неправильного вычислительных шагов в пункте «б» выделить нельзя.

\textbf{Почему это не является полным решением.} Выполненный пункт «а» не даёт ответа на пункт «б»: это два независимых вычисления. Для завершения задания необходимо определить число, куб которого равен $42{,}875$, и записать найденный результат.

\textcolor{blue!60!black}{\textbf{Ключевая идея:}} корень третьей степени из неотрицательного числа $a$ --- это число $b$, для которого $b^3=a$. Десятичную дробь можно представить обыкновенной дробью и извлечь корень из числителя и знаменателя, если они являются точными кубами.

\textbf{Исправление.} Не меняя решения пункта «а», дополним работу вычислением пункта «б»:
\[
42{,}875=\frac{42875}{1000}=\frac{343}{8}.
\]
\[
\sqrt[3]{42{,}875}=\sqrt[3]{\frac{343}{8}}=\frac{7}{2}=3{,}5.
\]

\textbf{Полное правильное решение.}

\textbf{Шаг 1.} Переведём десятичную дробь в обыкновенную. В числе $42{,}875$ три цифры после запятой, поэтому знаменатель первоначальной дроби равен $1000$:
\[
42{,}875=\frac{42875}{1000}.
\]

\textbf{Шаг 2.} Сократим числитель и знаменатель на $125$:
\[
\frac{42875}{1000}=\frac{343}{8}.
\]
Числа $343$ и $8$ являются кубами целых чисел:
\[
343=7^3,\qquad 8=2^3.
\]

\textbf{Шаг 3.} Используем свойство корня третьей степени из дроби:
\[
\sqrt[3]{\frac{343}{8}}=
\frac{\sqrt[3]{343}}{\sqrt[3]{8}}=
\frac{7}{2}.
\]

\textbf{Шаг 4.} Переведём обыкновенную дробь в десятичную:
\[
\frac72=3{,}5.
\]

\textbf{Проверка.} Возведём найденное число в куб:
\[
(3{,}5)^3=3{,}5\cdot3{,}5\cdot3{,}5
=12{,}25\cdot3{,}5=42{,}875.
\]
Результат совпадает с подкоренным числом, следовательно, вычисление верно.

\textcolor{teal!60!black}{\textbf{Итог:}} $\boxed{\sqrt[3]{42{,}875}=3{,}5}$.

\textbf{Задача на закрепление.} Вычислите $\sqrt[3]{15{,}625}$.

\newpage
\phantomsection\label{task:10}
\section*{Задание 10. Изменение уровня воды}
\textbf{Текст задания.} После дождя уровень воды в колодце может повыситься. Время $t$ падения небольшого камешка и расстояние $h$ от места его падения до воды связаны формулой $h=5t^2$, где $h$ измеряется в метрах, а $t$ --- в секундах. До дождя время падения было $0{,}6$ с. На сколько должен подняться уровень воды, чтобы время падения уменьшилось на $0{,}2$ с? Ответ выразите в метрах.

\textbf{Ваше решение.} В работе записано:
\[
\Delta h=5\cdot(0{,}2)^2=0{,}2\text{ м}.
\]

\textbf{Ваш ответ.} $0{,}2$ м.

\textcolor{red}{\textbf{Ошибка:}} изменение времени было ошибочно подставлено в формулу расстояния как полное время падения.

\textbf{Где именно возникает ошибка.} Верно, что расстояние до воды вычисляется по формуле $h=5t^2$, а время уменьшилось на $0{,}2$ с. Первый неверный переход --- запись
\[
\Delta h=5\cdot(\Delta t)^2.
\]
Эта запись не следует из исходной формулы. Число $0{,}2$ с описывает \emph{разницу между двумя временами}, а не время падения после дождя. Для определения повышения уровня воды надо отдельно найти старое и новое расстояния, а затем вычесть их.

\textbf{Почему этот шаг неверен.} Формула $h=5t^2$ устанавливает зависимость \emph{самого расстояния} от \emph{самого времени} падения. Если исходное время равно $t_1$, а новое равно $t_2$, то
\[
h_1=5t_1^2,\qquad h_2=5t_2^2.
\]
Поэтому изменение расстояния равно
\[
\Delta h=h_1-h_2=5(t_1^2-t_2^2),
\]
а не $5(t_1-t_2)^2$. В общем случае квадрат разности не равен разности квадратов:
\[
(t_1-t_2)^2=t_1^2-2t_1t_2+t_2^2,
\qquad t_1^2-t_2^2=(t_1-t_2)(t_1+t_2).
\]
Именно смешение этих двух выражений приводит к неправильному ответу.

\textcolor{blue!60!black}{\textbf{Ключевая идея:}} если величина зависит от времени по квадратичному закону, то для вычисления её изменения нужно сравнить значения функции \emph{в двух моментах времени}. Уменьшение времени само по себе нельзя подставлять вместо нового времени.

\textbf{Исправление от последнего правильного шага.} Формулу $h=5t^2$ сохраняем. Вместо подстановки только уменьшения времени сначала находим время после дождя:
\[
t_2=0{,}6-0{,}2=0{,}4\text{ с}.
\]
Теперь вычисляем разность расстояний:
\[
\Delta h=5\cdot(0{,}6)^2-5\cdot(0{,}4)^2=1\text{ м}.
\]

\textbf{Полное правильное решение.}

\textbf{Шаг 1.} Обозначим через $h_1$ расстояние от места падения камешка до поверхности воды до дождя. Тогда при $t_1=0{,}6$ с:
\[
h_1=5\cdot(0{,}6)^2.
\]
Квадрат числа $0{,}6$ равен $0{,}36$, поэтому
\[
h_1=5\cdot0{,}36=1{,}8\text{ м}.
\]

\textbf{Шаг 2.} Время после дождя на $0{,}2$ с меньше первоначального. Следовательно,
\[
t_2=t_1-0{,}2=0{,}6-0{,}2=0{,}4\text{ с}.
\]

\textbf{Шаг 3.} После дождя камешек долетает до воды за $0{,}4$ с. Новое расстояние $h_2$:
\[
h_2=5\cdot(0{,}4)^2=5\cdot0{,}16=0{,}8\text{ м}.
\]

\textbf{Шаг 4.} Уровень воды поднялся, поэтому расстояние от места падения камешка до воды \emph{уменьшилось}. Высота подъёма равна разности старого и нового расстояний:
\[
\Delta h=h_1-h_2=1{,}8-0{,}8=1\text{ м}.
\]

\textbf{Проверка.} После подъёма воды на $1$ м расстояние до поверхности составляет $1{,}8-1=0{,}8$ м. По исходной формуле
\[
5t^2=0{,}8,\qquad t^2=0{,}16.
\]
Время падения неотрицательно, значит $t=0{,}4$ с. Оно действительно на $0{,}2$ с меньше первоначальных $0{,}6$ с. Условие выполнено.

\textcolor{teal!60!black}{\textbf{Итог:}} $\boxed{1\text{ м}}$.

\textbf{Задача на закрепление.} Расстояние от точки падения камешка до поверхности воды определяется формулой $h=5t^2$. До подъёма уровня воды камешек падал $0{,}8$ с, а после подъёма --- на $0{,}3$ с меньше. На сколько метров поднялся уровень воды?

\vfill
\begin{center}
\small\itshape Лёвин Артём Александрович эксклюзивно для Григория Архипова
\end{center}
\end{document}